\documentclass[11pt]{article}
\usepackage{fontspec,amsmath,amssymb,amsthm,hyperref}
\usepackage[margin=1in]{geometry}
\setlength{\emergencystretch}{3em}
\title{Conditional entropy convergence for the actual Gaussian particle laws}
\author{LangevinLean mathematical exposition}
\date{}
\begin{document}\maketitle
\noindent Source module: \texttt{IPLA/LSI/Theorem.lean}. Compile this proof document with XeLaTeX or LuaLaTeX. The accompanying HTML places the complete Lean declarations beside every proof.


Let \(V_N(z)=\sum_iU(\theta,\sqrt Nz_i)\), \(D=d_\theta+Nd_x\), and \(\Pi_N\propto e^{-V_N}\). Assume differentiability and \(L\)-Lipschitz joint gradient, an LSI of constant \(\beta>0\) for this actual target, \(\beta\le NL\), and \(0<\gamma\le\beta/(4NL^2)\). Take an initial probability law \(\mu_0\). The endpoint also assumes \texttt{RegularInterpolationData} for \((V_N,\beta,\gamma/N,NL,\mu_0)\). These density, weak-equation, energy and moment obligations remain conditional; they have not been constructed from LSI and smoothness alone.


\section{The canonical particle laws are the actual Gaussian algorithm laws}


Define \(\rho_k=\operatorname{ulaLawStep}(V_N,\gamma/N)^{[k]}\mu_0\); at \(k=0\) this is \(\mu_0\) by the iterate definition. The lifted gradient is \(NL\)-Lipschitz, so the measurable Gaussian transition preserves probability, and induction through the shared ULA law theorem proves \(\rho_k\) is a probability for every \(k\). If \(\mu_0\) is the initial law of the explicit independent-Gaussian IPLA recursion, \texttt{laws\_eq\_ulaLawStep} identifies it with this same iterate. Thus the entropy result concerns the actual algorithm rather than a substituted reference process.



\noindent Lean declarations: \texttt{particleLaws}, \texttt{particleLaws\_zero}, \texttt{particleLaws\_isProbability}, \texttt{particleLaws\_eq\_laws}.

\section{Density regularity proves finite actual KL before real conversion}


The LSI assumption includes that \(\Pi_N\) is a probability measure. The interpolation data give positive relative densities \(r_k(t)\) with \(r_k(t)\Pi_N\) equal to the actual frozen-drift interpolation law, together with local entropy regularity and spatial smoothness. The shared finiteness theorem applies with \(h=\gamma/N\ge0\) and proves

\[\operatorname{KL}(\rho_k\Vert\Pi_N)<\infty\quad\text{for every }k.\]

The time-zero interpolation identity transfers this density-level KL finiteness to the actual iterate. In the probability-law application, the integrable entropy density also equals that actual finite KL. This guard precedes every use of \texttt{toReal}; infinite divergence is never turned into an artificial zero. The proof requires the supplied regularity, not a strong-convexity premise.



\noindent Lean declarations: \texttt{particle\_kl\_finite}.

\section{Derive entropy production and the one-step inequality}


Put \(h=\gamma/N\), \(\ell=NL\), and \(H_k=\operatorname{KL}(\rho_k\Vert\Pi_N)\). The data provide the weak interpolation equation, differentiation and integration regularity, an error field \(e_k(t)\), and the bounds

\[\int r_k(t)\|e_k(t)\|^2d\Pi_N\le t^2\ell^2G_k+2tD\ell^2,\qquad G_k\le\frac{4\ell^2}{\beta}H_k+2D\ell.\]

The shared calculus tests the weak equation against the logarithmic density and uses Young's inequality to derive \(H'(t)\le-3J(t)/4+\int r\|e\|^2d\Pi_N\). LSI gives \(J(t)\ge2\beta H(t)\). Since \(h\le\beta/(4\ell^2)\) and \(\beta\le\ell\), one has \(h\ell\le1/4\) and \(\beta h\le1/4\). The error bounds then imply

\[H'(t)\le-\frac32\beta H(t)+\frac\beta4H_k+\frac52hD\ell^2.\]

The shared affine differential comparison and \(e^{-3x/2}+x/4\le e^{-x}\) for \(0\le x\le1/4\) give \(H_{k+1}\le e^{-\beta h}H_k+6h^2D\ell^2\). The IPLA proof checks positivity of \(L\) from \(0<\beta\le NL\), translates the step restriction, and substitutes the exact lifted dimension and forcing. Therefore

\[H_{k+1}\le e^{-\beta\gamma/N}H_k+6\gamma^2DL^2.\]

Dividing by positive \(N\) gives the normalized recurrence with forcing \(6\gamma^2L^2(d_x+d_\theta/N)\). Neither recurrence is a field of the interpolation structure.



\noindent Lean declarations: \texttt{particle\_kl\_one\_step}, \texttt{particle\_kl\_one\_step\_normalized}.

\section{Iterate and simplify the finite-entropy bound}


Let \(q=e^{-\beta h}\). Repeated substitution in the one-step estimate gives

\[H_k\le q^kH_0+6h^2D\ell^2\sum_{j=0}^{k-1}q^j.\]

Because \(0<\beta h\le1/4\), the elementary bound \(1-e^{-\beta h}\ge3\beta h/4\) bounds the geometric sum by \(4/(3\beta h)\). Thus the shared theorem proves \(H_k\le e^{-\beta hk}H_0+8hD\ell^2/\beta\). Substitute \(h=\gamma/N\), \(\ell=NL\), and \(D=d_\theta+Nd_x\) to obtain

\[H_k\le e^{-\beta\gamma k/N}H_0+\frac{8\gamma ND L^2}{\beta}.\]

The proof invokes the shared bound after deriving every lifted numerical and regularity input, then uses the exact scaling identities. Dividing the entire inequality by positive \(N\) yields

\[\frac{H_k}{N}\le e^{-\beta\gamma k/N}\frac{H_0}{N}+\frac{8\gamma DL^2}{\beta}.\]

\noindent Lean declarations: \texttt{particle\_kl\_bound}, \texttt{particle\_kl\_bound\_normalized}.

\section{Specialize only when the lifted target supplies \ensuremath{\beta} = N\ensuremath{\alpha}}


Assume the actual lifted target has the supplied LSI constant \(N\alpha\), with \(0<\alpha\le L\), and \(0<\gamma\le\alpha/(4L^2)\). Then \(N\alpha>0\), compatibility \(N\alpha\le NL\) holds, and the step equivalence supplies the general theorem's restriction. Its rate cancels to \(e^{-\alpha\gamma k}\); the normalized bias uses \(D/N=d_x+d_\theta/N\). Consequently

\[\frac{H_k}{N}\le e^{-\alpha\gamma k}\frac{H_0}{N}+\frac{8\gamma L^2}{\alpha}\left(d_x+\frac{d_\theta}{N}\right).\]

This is a specialization of a lifted-target LSI already given in the hypotheses. The shared parameter makes the lifted model nonproduct, so no pair-target LSI tensorization is inferred.



\noindent Lean declarations: \texttt{particle\_kl\_bound\_of\_linear\_lsi}.

\section{Transfer the canonical bounds to the supplied Gaussian implementation}


For a square-integrable initialization on \(P\) and noises satisfying all scalar Gaussian and independence fields, the law identity gives \(\rho_k=\operatorname{laws}(P,U,\gamma,k)\). Rewriting the normalized general bound by this equality proves the first Gaussian implementation corollary. Rewriting the specialized \(N\alpha\)-LSI bound proves the second one, with all algorithm and interpolation assumptions retained. The entropy on the right is the full initial joint divergence relative to \(\Pi_N\); normalization divides both it and the subsequent joint divergence by \(N\). These conclusions by themselves do not assert parameter concentration about a mode.



\noindent Lean declarations: \texttt{gaussian\_ipla\_kl\_bound\_normalized}, \texttt{gaussian\_ipla\_kl\_bound\_of\_linear\_lsi}.

\end{document}
